\documentclass[12pt]{article}

% ------------------------------------------------
\usepackage[spanish]{babel}
\usepackage[utf8]{inputenc}
\usepackage[T1]{fontenc}
\usepackage{geometry}
\geometry{margin=2.5cm}
\usepackage{setspace}
\onehalfspacing
\usepackage{parskip}
\usepackage[hypertexnames=false]{hyperref}
\usepackage{fancyhdr}
\usepackage{siunitx}
\sisetup{per-mode=symbol}
\usepackage{float}
\usepackage{graphicx}
\usepackage{adjustbox}
\usepackage{tikz}
\usetikzlibrary{arrows.meta,positioning,shapes.geometric}

\pagestyle{fancy}
\fancyhf{}
\rhead{\thepage}
\setlength{\headheight}{14.5pt}

\begin{document}

% ------------------------------------------------
\begin{center}
\textbf{\Large Universidad de Costa Rica}\\
\vspace{0.3cm}
\textbf{Escuela de Ingeniería Eléctrica}

\vspace{2cm}

\textbf{\LARGE Bitácora de laboratorio}

\vspace{1.5cm}

\textbf{Diseño de PCB para Sensores de Torque y Temperatura}

\vspace{1.5cm}

Kenneth Daniel Vizcaíno Jiménez -- C09331

\vspace{1cm}

Marzo--Junio 2026
\end{center}

\newpage
\tableofcontents
\newpage

% =================================================
\section{Descripción del proyecto}

\textbf{Título del proyecto:}
Diseño e implementación de una tarjeta de acondicionamiento de señales para sensores de torque y temperatura en un motor de inducción.

\textbf{Objetivo general:}
Diseñar e implementar una PCB que permita acondicionar las señales de sensores para su correcta adquisición.

\textbf{Objetivos específicos:}
\begin{enumerate}
\item Comprender el diseño de PCB.
\item Verificar sensores.
\item Diseñar circuito.
\item Implementar prototipo.
\end{enumerate}

\textbf{Entregables:}
\begin{itemize}
\item Esquemático
\item PCB
\item Prototipo
\item Pruebas
\end{itemize}

% =================================================
\section{Bitácora de trabajo}

% =========================
\section{Primera semana}

\subsection{Inicio del proyecto}
\textbf{Fecha:} 11/03/2026 \quad \textbf{Hora:} 13:00--14:00

\textbf{Objetivo:}
Definir el enfoque del proyecto.

\textbf{Desarrollo:}
\begin{itemize}
\item Se realizó una revisión general del alcance del proyecto, identificando los requerimientos principales del sistema y los objetivos a cumplir en el diseño de la PCB.
\item Se planteó una primera idea conceptual del circuito, considerando la integración simultánea de sensores de temperatura y torque dentro de una misma tarjeta.
\item Se analizaron posibles configuraciones de conexión eléctrica, evaluando cómo se distribuirían las señales dentro del sistema.
\end{itemize}

\textbf{Resultados:}
Definición inicial del sistema.

\textbf{Dificultades:}
Incertidumbre sobre componentes.

% -------------------------
\subsection{Introducción a KiCad}
\textbf{Fecha:} 12/03/2026 \quad \textbf{Hora:} 13:00--17:00

\textbf{Objetivo:}
Aprender el entorno de KiCad.

\textbf{Desarrollo:}
\begin{itemize}
\item Se estudiaron tutoriales enfocados en el uso básico del software, permitiendo comprender la lógica de trabajo del entorno.
\item Se exploraron los conceptos de símbolos electrónicos y footprints, entendiendo su relación dentro del diseño.
\item Se revisaron distintas librerías de componentes disponibles para su implementación en el proyecto.
\end{itemize}

\textbf{Resultados:}
Comprensión funcional del software.

\textbf{Dificultades:}
Asociación símbolo-huella.

% -------------------------
\subsection{Práctica de PCB}
\textbf{Fecha:} 13/03/2026 \quad \textbf{Hora:} 11:00--16:00

\textbf{Objetivo:}
Practicar diseño.

\textbf{Desarrollo:}
\begin{itemize}
\item Se desarrolló un prototipo de PCB con fines de aprendizaje, sin ser aún el diseño final del proyecto.
\item Se trabajó en la colocación de componentes y su distribución dentro del área de la tarjeta.
\item Se asignaron footprints correspondientes a cada componente, fortaleciendo el flujo de trabajo.
\end{itemize}

\textbf{Resultados:}
Mayor dominio práctico.

\textbf{Dificultades:}
Orden del diseño.

% =========================
\section{Segunda semana}

\subsection{Profundización en diseño}
\textbf{Fecha:} 18/03/2026 \quad \textbf{Hora:} 15:00--23:00

\textbf{Objetivo:}
Mejorar diseño.

\textbf{Desarrollo:}
\begin{itemize}
\item Se trabajó en el diseño del plano de tierra, comprendiendo su importancia en la reducción de ruido.
\item Se ajustó el contorno físico de la PCB, definiendo su forma final.
\item Se importaron nuevas librerías para ampliar las opciones de diseño.
\item Se revisaron reglas eléctricas y de manufactura para evitar errores futuros.
\end{itemize}

\textbf{Resultados:}
Mayor comprensión global.

\textbf{Dificultades:}
Implementación de tierra.

% =========================
\section{Tercera semana}

\subsection{Ruido en sistemas}
\textbf{Fecha:} 25/03/2026 \quad \textbf{Hora:} 13:30--18:30

\textbf{Objetivo:}
Comprender ruido.

\textbf{Desarrollo:}
\begin{itemize}
\item Se analizaron diferentes fuentes de ruido presentes en sistemas electrónicos.
\item Se estudiaron efectos térmicos e interferencias externas.
\item Se evaluaron técnicas como filtrado, blindaje y cableado trenzado para mitigación.
\end{itemize}

\textbf{Resultados:}
Identificación de soluciones prácticas.

\textbf{Dificultades:}
Conceptos abstractos.

% -------------------------
\subsection{Sensores}
\textbf{Fecha:} 25/03/2026 \quad \textbf{Hora:} 13:30--17:30

\textbf{Objetivo:}
Analizar sensores.

\textbf{Desarrollo:}
\begin{itemize}
\item Se investigó el funcionamiento de sensores de torque y temperatura.
\item Se analizaron sus señales de salida y características eléctricas.
\item Se evaluaron requisitos de acondicionamiento para integración futura.
\end{itemize}

\textbf{Resultados:}
Comprensión del sistema.

\textbf{Dificultades:}
Rangos de señal.

% =========================
\section{Cuarta semana}

\subsection{Revisión técnica}
\textbf{Fecha:} 30/03/2026 \quad \textbf{Hora:} 13:30--19:30

\textbf{Objetivo:}
Refuerzo teórico.

\textbf{Desarrollo:}
\begin{itemize}
\item Se analizaron documentos técnicos relacionados con diseño de PCB.
\item Se estudiaron buenas prácticas de ingeniería.
\item Se relacionó la teoría con el proyecto en desarrollo.
\end{itemize}

\textbf{Resultados:}
Mejor criterio técnico.

\textbf{Dificultades:}
Interpretación normativa.

% =========================
\section{Quinta semana}

% -------------------------
\subsection{Desarrollo de arquitectura del sistema}
\textbf{Fecha:} 08/04/2026 \quad \textbf{Hora:} 13:00--20:00

\textbf{Objetivo:}
Definir formalmente la arquitectura completa del sistema de adquisición de señales para los sensores de torque y temperatura.

\textbf{Desarrollo:}
\begin{itemize}
\item Se elaboró un documento técnico completo donde se definió la arquitectura del sistema, estableciendo la cadena funcional desde los sensores hasta el controlador BoomBox.
\item Se planteó un diagrama de bloques estructurado considerando las etapas: sensor, transporte por cable RJ45/UTP, entrada a PCB, acondicionamiento analógico, interfaz ModuLink y digitalización en el sistema imperix.
\item Se analizó en detalle el comportamiento del sensor de torque basado en puente de Wheatstone, identificando que su salida es diferencial y de muy baja amplitud (del orden de milivoltios), lo que hace indispensable el uso de amplificación instrumental.
\item Se establecieron valores iniciales de diseño para el sensor de torque, considerando sensibilidades típicas (2 mV/V a 3 mV/V) y estimando salidas de aproximadamente 10 mV a fondo de escala con excitación de 5 V.
\item Se definió la necesidad de implementar una etapa de acondicionamiento analógico que incluya amplificación, filtrado, offset y escalado antes de la digitalización.
\item Se analizó el sensor de temperatura real del motor, identificando que corresponde a un sistema PTC en serie utilizado para protección térmica, cuya salida es resistiva y no adecuada para medición continua.
\item Se propuso como mejora futura la incorporación de un sensor tipo Pt100, permitiendo una medición de temperatura más precisa mediante excitación de corriente y lectura diferencial.
\item Se evaluó el uso de cable RJ45/UTP como medio físico de interconexión, destacando el uso de pares trenzados para señales diferenciales y su impacto positivo en la inmunidad al ruido.
\item Se definieron criterios de diseño para la entrada de la PCB, incluyendo protección contra ESD, filtrado EMI y referencia de tierra adecuada.
\item Se dimensionó preliminarmente la ganancia del sistema de torque, estableciendo rangos de 500 a 1000 V/V para adaptar señales de milivoltios a niveles compatibles con el controlador.
\item Se identificaron integrados recomendables para el acondicionamiento, como amplificadores de instrumentación tipo INA333, INA826 y AD8426.
\item Se analizó la etapa de adquisición del sistema imperix BoomBox, comprendiendo su funcionamiento con ADC diferencial de 16 bits y rango de entrada de hasta ±10 V.
\item Se estableció la importancia de diseñar la PCB de forma que la señal final quede dentro del rango del ADC, evitando saturación o activación de protecciones.
\end{itemize}

\textbf{Resultados:}
Se obtuvo una arquitectura completa y fundamentada del sistema, sirviendo como base sólida para el diseño de la PCB y la integración con el controlador.

\textbf{Dificultades:}
Integrar correctamente todos los bloques del sistema y validar supuestos de diseño sin contar aún con todos los datasheets definitivos.

% =========================
\subsection{Validación experimental de sensores de temperatura}
\textbf{Fecha:} 15/04/2026 \quad \textbf{Hora:} 14:00--18:30

\textbf{Objetivo:}
Verificar experimentalmente el tipo de sensor de temperatura presente en el motor mediante una visita al laboratorio LabCES.

\textbf{Desarrollo:}
\begin{itemize}
\item Se realizó una visita al laboratorio \textbf{LabCES} con el fin de inspeccionar directamente el motor y verificar la naturaleza de los sensores de temperatura instalados.
\item Se efectuó una inspección física del cableado del motor, identificando la presencia de tres grupos de conductores similares, cada uno compuesto por tres hilos.
\item Se planteó la hipótesis de que cada grupo correspondía a un sensor independiente asociado a cada fase del motor trifásico.
\item Se procedió a realizar mediciones de resistencia utilizando un multímetro en cada uno de los grupos.
\item Se obtuvieron valores en el rango de \SIrange{111}{112}{\ohm} de manera consistente en los tres conjuntos de conductores.
\item A partir de estos resultados, se descartó la hipótesis inicial de sensores tipo PTC y se identificó que los sensores corresponden a \textbf{RTD tipo Pt100 de 3 hilos}.
\item Se relacionaron los valores medidos con el comportamiento teórico del Pt100, el cual presenta una resistencia nominal de \SI{100}{\ohm} a \SI{0}{\celsius} y aumenta con la temperatura.
\end{itemize}

\textbf{Resultados:}
Se confirmó experimentalmente, durante la visita a LabCES, que el sistema de temperatura del motor está compuesto por \textbf{tres sensores Pt100 de 3 hilos}, uno por fase.

\textbf{Dificultades:}
Interpretación inicial del cableado y diferenciación entre sensores PTC y RTD.

% -------------------------

\subsection{Rediseño del sistema y ajuste de parámetros de adquisición}
\textbf{Fecha:} 17/04/2026 \quad \textbf{Hora:} 19:00--23:30

\textbf{Objetivo:}
Actualizar la arquitectura del sistema y definir parámetros de ganancia y margen de seguridad.

\textbf{Desarrollo:}
\begin{itemize}
\item Se actualizó el documento técnico del proyecto, incorporando la validación experimental de los sensores Pt100 y eliminando la referencia previa a sensores PTC.
\item Se redefinió la arquitectura del sistema, estableciendo tres canales independientes de temperatura y un canal de torque.
\item Se analizó el rango de salida del sistema de referencia imperix, identificando que el módulo ModuLink presenta una salida máxima aproximada de \(\pm 4.95\ \mathrm{V}\).
\item Se tomó la decisión de no operar al límite del sistema, sino de implementar un margen de seguridad del orden del \(5\%\), definiendo una tensión eléctrica máxima de diseño cercana a \SI{4.7}{V}.
\item Para el sensor de torque, se adoptó un valor de diseño intermedio de \SI{50}{mV}, dentro del rango típico de sensores tipo puente, y se calculó una ganancia aproximada de \(95\) a \(100\ \mathrm{V/V}\).
\item Para los sensores Pt100, considerando una señal aproximada de \SI{112}{mV} obtenida mediante excitación por corriente, se estimó una ganancia de \(40\) a \(42\ \mathrm{V/V}\).
\item Se reforzó la necesidad de amplificación mediante amplificadores de instrumentación para el canal de torque, debido a su baja amplitud y naturaleza diferencial.
\item Se estableció que el canal de temperatura requiere una etapa previa de conversión de resistencia a tensión eléctrica antes del acondicionamiento.
\item Se ajustaron las ecuaciones del sistema, los diagramas de bloques y la tabla de señales, logrando coherencia entre el diseño teórico y los resultados experimentales.
\end{itemize}

% ---- AQUÍ VA EL DIAGRAMA ----

\begin{figure}[H]
\centering
\begin{adjustbox}{max width=0.98\textwidth}
\begin{tikzpicture}[
node distance=0.65cm and 0.85cm,
block/.style={draw, rounded corners, align=center, minimum height=1.15cm, minimum width=2.8cm},
line/.style={-{Latex[length=2.2mm]}}
]

\node[block] (torque) {Sensor de torque\\(puente de Wheatstone)};
\node[block, below=of torque] (temp) {3 sensores Pt100\\(uno por fase)};

\node[block, right=of torque] (in1) {Protección\\y filtrado};
\node[block, right=of in1] (amp1) {Amplificación\\instrumental};
\node[block, right=of amp1] (drv1) {Driver\\diferencial};

\node[block, right=of temp] (r2v) {Excitación\\y conversión R$\rightarrow$V};
\node[block, right=of r2v] (amp2) {Amplificación\\y filtrado};
\node[block, right=of amp2] (drv2) {Driver\\diferencial};

\node[block, right=1.3cm of drv1, minimum width=3.6cm] (stp) {Cable STP / RJ45\\blindado};
\node[block, right=of stp, minimum width=3.7cm] (bbox) {BoomBox\\ADC diferencial};
\node[block, below=of bbox, minimum width=3.0cm] (ctrl) {FPGA / CPU\\control y monitoreo};

\draw[line] (torque) -- (in1);
\draw[line] (in1) -- (amp1);
\draw[line] (amp1) -- (drv1);
\draw[line] (temp) -- (r2v);
\draw[line] (r2v) -- (amp2);
\draw[line] (amp2) -- (drv2);

\draw[line] (drv1.east) -- ++(0.35,0) |- (stp.west);
\draw[line] (drv2.east) -- ++(0.35,0) |- (stp.west);

\draw[line] (stp) -- (bbox);
\draw[line] (bbox) -- (ctrl);

\end{tikzpicture}
\end{adjustbox}
\caption{Diagrama de bloques actualizado del sistema de adquisición, incorporando los tres sensores Pt100 y el acondicionamiento diferencial hacia BoomBox.}
\label{fig:bloques_semana6}
\end{figure}

% ---- FIN DIAGRAMA ----

\textbf{Resultados:}
Se obtuvo una versión actualizada y consistente del diseño, con parámetros de ganancia y márgenes de seguridad definidos y justificados.

\textbf{Dificultades:}
Definir valores de diseño balanceando resolución, ruido y saturación sin exceder el rango del sistema.

% =========================
\section{Sexta semana}

% -------------------------
\subsection{Análisis del sistema de adquisición y arquitectura general}
\textbf{Fecha:} 22/04/2026 \quad \textbf{Hora:} 18:00--20:30

\textbf{Objetivo:}
Comprender la arquitectura completa del sistema y definir el flujo de señal desde los sensores hasta la BoomBox.

\textbf{Desarrollo:}
\begin{itemize}
\item Se revisó la arquitectura general del sistema, identificando que la señal debe mantenerse analógica hasta la BoomBox.
\item Se corrigió la idea inicial de digitalización en la PCB, estableciendo que el ADC pertenece al sistema imperix.
\item Se definió el uso de señal diferencial como requisito fundamental para la transmisión.
\item Se estableció la cadena funcional sensor--PCB--BoomBox como base del diseño.
\end{itemize}

\textbf{Resultados:}
Se definió correctamente la arquitectura del sistema de adquisición.

\textbf{Dificultades:}
Diferenciar claramente entre acondicionamiento analógico y digitalización.

% -------------------------
\subsection{Análisis de sensores y cálculo de estimación de parámetros }
\textbf{Fecha:} 23/04/2026 \quad \textbf{Hora:} 19:00--22:00

\textbf{Objetivo:}
Analizar cómo se comportan eléctricamente los sensores y definir los valores básicos necesarios para el diseño.

\textbf{Desarrollo:}
\begin{itemize}
\item Se analizó el sensor de torque como un puente de Wheatstone, identificando su salida en el rango de milivolts.
\item Se estableció un valor de diseño de \SI{50}{mV} para el cálculo de ganancia.
\item Se analizaron los sensores Pt100, confirmando su comportamiento resistivo y la necesidad de excitación por corriente.
\item Se definió una corriente de excitación de \SI{1}{mA} para los sensores de temperatura.
\item Se calcularon las señales esperadas en los sensores, obteniendo valores del orden de \SIrange{100}{112}{mV}.
\end{itemize}

\textbf{Resultados:}
Se definieron los rangos eléctricos de entrada del sistema.

\textbf{Dificultades:}
Ajustar valores sin contar aún con todos los datasheets definitivos.

% -------------------------
\subsection{Definición de etapas de amplificación y filtrado}
\textbf{Fecha:} 23/04/2026 \quad \textbf{Hora:} 22:00--23:30

\textbf{Objetivo:}
Definir las etapas de amplificación y la estrategia de filtrado del sistema.

\textbf{Desarrollo:}
\begin{itemize}
\item Se estableció el uso de amplificadores de instrumentación para la lectura del sensor de torque.
\item Se definieron ganancias aproximadas de \(G \approx 100\) para torque y \(G \approx 40\) para temperatura.
\item Se analizó el uso de amplificadores operacionales auxiliares para etapas de filtrado y ajuste.
\item Se propuso la implementación de un filtro pasabajo ajustable mediante trimmer.
\item Se definió un rango preliminar de frecuencia de corte para validación experimental.
\end{itemize}

\textbf{Resultados:}
Se establecieron las bases del acondicionamiento analógico.

\textbf{Dificultades:}
Determinar la frecuencia de corte sin mediciones reales de ruido.

% -------------------------
\subsection{Definición final del esquema de acondicionamiento y bloques de la PCB}
\textbf{Fecha:} 24/04/2026 \quad \textbf{Hora:} 17:00--23:00

\textbf{Objetivo:}
Consolidar el diseño del esquema eléctrico preliminar de la PCB, integrando todos los bloques funcionales, valores de diseño y criterios de ingeniería definidos durante la semana.

\textbf{Desarrollo:}
\begin{itemize}
\item Se realizó una revisión completa del sistema de acondicionamiento, integrando la información obtenida sobre los sensores de torque y temperatura en una única arquitectura coherente.
\item Se definió formalmente la cadena de procesamiento de señal desde los sensores hasta la BoomBox, estableciendo las etapas: protección, amplificación instrumental, filtrado, ajuste de escala y salida diferencial.
\item Se adoptó un valor de diseño de \SI{50}{mV} para el sensor de torque y se definió una ganancia aproximada de \(G \approx 100\).
\item Para los sensores Pt100, se estableció una corriente de excitación de \SI{1}{mA} y una ganancia aproximada de \(G \approx 40\).
\item Se definió la inclusión de un filtro pasabajo ajustable mediante trimmer para validar experimentalmente el ruido del sistema.
\item Se estableció el uso de señal diferencial en la salida de la PCB y cable STP con conector RJ45 como medio físico.
\item Se definieron criterios de alimentación mediante fuente externa \(V_{dc}\), con consumo estimado entre \SIrange{30}{70}{mA}.
\item Se incorporaron elementos de protección y desacoplo para robustecer la PCB.
\end{itemize}

\textbf{Resultados:}
Se obtuvo una definición clara, coherente y fundamentada del esquema eléctrico de la PCB.

\textbf{Dificultades:}
Definir valores de diseño sin mediciones completas del entorno real.

\section{Séptima semana}

% -------------------------
\subsection{Validación de parámetros eléctricos mediante hojas de datos}
\textbf{Fecha:} 29/04/2026 \quad \textbf{Hora:} 18:00--22:00

\textbf{Objetivo:}
Verificar los valores eléctricos reales de los sensores y ajustar los parámetros de diseño con base en información técnica confiable.

\textbf{Desarrollo:}
\begin{itemize}
\item Se revisaron hojas de datos del sensor de torque tipo puente Wheatstone, confirmando una impedancia de entrada y salida cercana a \SI{350}{\ohm}.
\item Se validó una excitación recomendada de \SI{10}{V}, con una sensibilidad aproximada de \SI{3}{mV/V}.
\item Se determinó que la salida máxima del sensor es del orden de \SIrange{1}{30}{mV}, consistente con sensores de este tipo \cite{InterfaceSM50}.
\item Se analizaron las características de precisión del sensor, observando bajo error de cero, baja no linealidad y baja histéresis.
\item Se confirmó que los sensores de temperatura corresponden a tipo Pt100, basado en mediciones experimentales realizadas en laboratorio.
\item Se verificó el comportamiento resistivo del Pt100, con valores de \SI{100}{\ohm} a \SI{0}{^\circ C} y aproximadamente \SI{138.5}{\ohm} a \SI{100}{^\circ C}.
\item Se validó el uso de una corriente de excitación de \SI{1}{mA}, obteniendo señales en el rango de \SIrange{100}{138}{mV}.
\end{itemize}

\textbf{Resultados:}
Se obtuvieron valores eléctricos reales y confiables para ambos sensores, permitiendo validar el diseño del sistema.

\textbf{Dificultades:}
Falta de identificación directa del sensor de temperatura, lo que requirió validación experimental.

% -------------------------
\subsection{Análisis del acondicionamiento mediante puente Wheatstone}
\textbf{Fecha:} 29/04/2026 \quad \textbf{Hora:} 22:00--23:30

\textbf{Objetivo:}
Definir la estrategia de obtención de señal para sensores resistivos y su integración en el sistema.

\textbf{Desarrollo:}
\begin{itemize}
\item Se analizó el funcionamiento del puente Wheatstone como método de medición diferencial para sensores resistivos.
\item Se estableció que el sensor de torque ya integra un puente completo, mientras que el Pt100 requiere configuración externa.
\item Se definió que la señal útil es la diferencia entre las salidas del puente, la cual es proporcional a la variación de resistencia.
\item Se identificó que las señales obtenidas son del orden de \SIrange{1}{30}{mV}, requiriendo amplificación.
\item Se reafirmó el uso de amplificadores de instrumentación debido a su alta impedancia de entrada y rechazo de modo común .
\end{itemize}

\textbf{Resultados:}
Se definió correctamente la forma de obtención de señal para ambos sensores.

\textbf{Dificultades:}
Comprender el impacto del modo común en la medición diferencial.

% -------------------------
\subsection{Revisión con el profesor y validación de salidas del sistema}
\textbf{Fecha:} 30/04/2026 \quad \textbf{Hora:} 20:00--23:00

\textbf{Objetivo:}
Validar con el profesor los niveles de salida del sistema y confirmar la viabilidad del diseño.

\textbf{Desarrollo:}
\begin{itemize}
\item Se presentó al profesor el análisis realizado de los sensores y sus parámetros eléctricos.
\item Se discutieron los niveles de señal obtenidos y su compatibilidad con la etapa de adquisición (BoomBox).
\item Se validaron formalmente los niveles de salida definidos, respaldados por hojas de datos.
\item Se recibió retroalimentación sobre la importancia del manejo diferencial de la señal.
\item Se discutieron criterios de robustez frente a ruido en el entorno industrial del motor.
\item El profesor proporcionó material bibliográfico sobre diseño de sistemas de acondicionamiento de señal para fortalecer la base teórica del proyecto.
\end{itemize}

\textbf{Resultados:}
Se validaron los niveles de salida del sistema y se confirmó la coherencia del diseño propuesto.

\textbf{Dificultades:}
Ajustar el diseño considerando restricciones reales del sistema de adquisición.

% -------------------------
\subsection{Integración final de parámetros y ajustes del diseño}
\textbf{Fecha:} 30/04/2026 \quad \textbf{Hora:} 17:00--18:30

\textbf{Objetivo:}
Integrar todos los parámetros validados en el diseño del sistema de acondicionamiento.

\textbf{Desarrollo:}
\begin{itemize}
\item Se ajustaron los valores de ganancia en función de los datos reales obtenidos.
\item Se consolidó el uso de \(G \approx 100\) para torque y \(G \approx 40\) para temperatura.
\item Se verificó la coherencia entre los niveles de entrada y salida del sistema.
\item Se integraron los criterios de diseño en la arquitectura general de la PCB.
\item Se dejaron definidos los parámetros base para la siguiente etapa de diseño esquemático en KiCad.
\end{itemize}

\textbf{Resultados:}
Se obtuvo un diseño ajustado y validado con base en datos reales y criterio técnico.

\textbf{Dificultades:}
Ajustar todos los parámetros manteniendo compatibilidad con el sistema completo.


% =========================
\section{Octava semana}

% -------------------------
\subsection{Revisión de avances y organización de la bitácora}
\textbf{Fecha:} 05/05/2026 \quad \textbf{Hora:} 18:00--21:00

\textbf{Objetivo:}
Ordenar los avances acumulados del proyecto y dejar documentadas las decisiones técnicas tomadas hasta el momento.

\textbf{Desarrollo:}
\begin{itemize}
\item Se revisó el avance general de la bitácora para verificar que las semanas anteriores mantuvieran una estructura uniforme.
\item Se reorganizó la información relacionada con la arquitectura del sistema de acondicionamiento.
\item Se reforzó la idea de mantener la señal analógica desde los sensores hasta la etapa de adquisición del sistema BoomBox.
\item Se revisaron nuevamente los bloques principales del sistema: sensor, protección, amplificación, filtrado, ajuste de escala y salida diferencial.
\item Se corrigió la redacción técnica del documento para usar términos más adecuados, como tensión eléctrica, volts y milivolts.
\item Se identificaron puntos que aún requerían justificación con hojas de datos y material técnico.
\end{itemize}

\textbf{Resultados:}
Se dejó más ordenada la documentación del proyecto y se aclararon las decisiones principales de la arquitectura de la PCB.

\textbf{Dificultades:}
Mantener la bitácora clara sin perder información técnica importante del proyecto.

% -------------------------
\subsection{Revisión de criterios para el acondicionamiento de señales}
\textbf{Fecha:} 10/05/2026 \quad \textbf{Hora:} 19:00--22:30

\textbf{Objetivo:}
Reforzar la base teórica necesaria para el diseño de la PCB mediante la revisión de conceptos de sensores y acondicionamiento de señal.

\textbf{Desarrollo:}
\begin{itemize}
\item Se inició una revisión guiada del libro de sensores y acondicionadores de señal, enfocándose en los temas más relacionados con el proyecto.
\item Se identificaron conceptos importantes como sensibilidad, rango dinámico, ganancia, offset, resolución e incertidumbre.
\item Se relacionaron estos conceptos con las señales esperadas en el sensor de torque y en los sensores de temperatura.
\item Se analizó la necesidad de adaptar señales pequeñas, del orden de milivolts, hacia niveles compatibles con la adquisición.
\item Se repasó la importancia de evitar saturación en las etapas analógicas y de dejar margen suficiente en el rango de salida.
\item Se anotaron fórmulas base para continuar con los cálculos de ganancia y escalamiento durante la siguiente semana.
\end{itemize}

\textbf{Resultados:}
Se reforzó la teoría necesaria para justificar el diseño del acondicionamiento y se identificaron las fórmulas principales para el análisis de señales.

\textbf{Dificultades:}
Distinguir cuándo una tensión corresponde al valor mínimo, máximo, offset o fondo de escala del sistema.

% =========================
\section{Novena semana}

% -------------------------
\subsection{Estudio de rango dinámico, offset y resolución del ADC}
\textbf{Fecha:} 13/05/2026 \quad \textbf{Hora:} 16:00--19:00

\textbf{Objetivo:}
Comprender cómo se relacionan la señal del sensor, la ganancia, el offset y la resolución del ADC dentro del sistema de adquisición.

\textbf{Desarrollo:}
\begin{itemize}
\item Se estudiaron problemas del libro relacionados con rango dinámico, resolución y acondicionamiento de señales.
\item Se revisó el significado de términos como fondo de escala, valor mínimo, valor máximo, offset y número de bits.
\item Se analizó cómo una señal pequeña del sensor puede ocupar solo una parte del rango del ADC si no se amplifica correctamente.
\item Se revisó la expresión del efecto del offset en una cadena de medición, considerando resistencia de salida, corriente de entrada y ganancia.
\item Se estudió cómo una tensión de offset amplificada puede reducir el número efectivo de bits disponibles para representar la señal útil.
\item Se relacionó este análisis con el diseño del acondicionamiento para torque y temperatura.
\end{itemize}

\textbf{Resultados:}
Se comprendió mejor la función del offset y su impacto sobre la resolución efectiva del sistema de adquisición.

\textbf{Dificultades:}
Interpretar correctamente si una tensión dada representa el valor mínimo, el offset o el valor máximo esperado del sensor.

% -------------------------
\subsection{Análisis de fórmulas para puente Wheatstone y señal diferencial}
\textbf{Fecha:} 14/05/2026 \quad \textbf{Hora:} 18:00--21:30

\textbf{Objetivo:}
Revisar las ecuaciones necesarias para analizar un puente Wheatstone aplicado a sensores resistivos.

\textbf{Desarrollo:}
\begin{itemize}
\item Se repasó el funcionamiento general del puente Wheatstone como estructura de medición para sensores resistivos.
\item Se revisó que, cuando el puente está balanceado, la salida diferencial ideal es cercana a cero.
\item Se analizó que una variación de resistencia en los brazos del puente produce una tensión diferencial proporcional al cambio físico medido.
\item Se relacionó el puente Wheatstone con el sensor de torque, debido a que este tipo de sensor entrega una señal diferencial de baja amplitud.
\item Se reforzó la necesidad de usar un amplificador de instrumentación para medir la diferencia entre las salidas del puente y rechazar la tensión de modo común.
\item Se revisó la importancia de que la etapa de entrada tenga alta impedancia para no cargar el puente ni alterar la medición.
\end{itemize}

\textbf{Resultados:}
Se aclaró la relación entre puente Wheatstone, señal diferencial, modo común y amplificación instrumental.

\textbf{Dificultades:}
Diferenciar entre la tensión diferencial útil del puente y la tensión de modo común presente en sus nodos de salida.

% -------------------------
\subsection{Consolidación de fórmulas y criterios para el diseño de la PCB}
\textbf{Fecha:} 15/05/2026 \quad \textbf{Hora:} 17:00--20:00

\textbf{Objetivo:}
Consolidar las fórmulas y criterios técnicos necesarios para continuar con el diseño esquemático de la PCB.

\textbf{Desarrollo:}
\begin{itemize}
\item Se organizaron las fórmulas principales relacionadas con sensibilidad, ganancia, offset, rango dinámico y resolución del ADC.
\item Se revisó cómo nombrar correctamente las variables para evitar confusiones al aplicar las ecuaciones.
\item Se relacionaron las fórmulas del libro con el caso práctico del proyecto de sensores de torque y temperatura.
\item Se identificó que el diseño debe mantener margen de salida para evitar saturación frente a offset, ruido o variaciones reales del sensor.
\item Se reforzó la idea de que el acondicionamiento debe incluir protección, amplificación, filtrado y salida diferencial.
\item Se dejaron definidos los puntos teóricos que deben revisarse antes de avanzar con el esquemático final en KiCad.
\end{itemize}

\textbf{Resultados:}
Se obtuvo una base de fórmulas más clara para continuar el diseño de la PCB con mejor criterio técnico.

\textbf{Dificultades:}
Ordenar las fórmulas de forma que sean útiles para resolver problemas y, al mismo tiempo, aplicables al diseño real del proyecto.



% =========================
\section{Décima semana}

% -------------------------
\subsection{Validación de datos reales del sensor SM-50}
\textbf{Fecha:} 19/05/2026 \quad \textbf{Hora:} 10:00--14:00

\textbf{Objetivo:}
Corregir los datos del sensor de torque a partir de la información real de la hoja técnica del SM-50.

\textbf{Desarrollo:}
\begin{itemize}
\item Se revisó nuevamente la información del sensor Interface SM-50 para evitar trabajar con valores genéricos o supuestos incorrectos.
\item Se corrigió la capacidad utilizada en el proyecto, tomando como referencia el valor real de \SI{20}{in.lb}, equivalente aproximadamente a \SI{2.2}{N.m}.
\item Se confirmó que la salida nominal del sensor es de \SI{3.0}{mV/V}, por lo que con una excitación de \SI{10}{V} se obtiene una señal diferencial máxima cercana a \SI{30}{mV}.
\item Se registraron parámetros importantes para la defensa técnica: resistencia nominal del puente de \SI{350}{\ohm}, no linealidad de \SI{\pm 0.03}{\percent} FS, histéresis de \SI{\pm 0.02}{\percent} FS y balance de cero de \SI{\pm 1}{\percent} RO.
\item Se corrigió la interpretación del cableado del sensor, dejando como referencia: rojo para +EXC, negro para -EXC, verde para +OUT, blanco para -OUT y blindaje para SHIELD.
\end{itemize}

\textbf{Resultados:}
Se obtuvo una base de datos real para el canal de torque, eliminando el error de trabajar con rangos de torque que no correspondían al sensor seleccionado.

\textbf{Dificultades:}
Distinguir entre valores típicos de sensores de puente y los valores específicos del SM-50 que realmente se usarán en el proyecto.

% -------------------------
\subsection{Corrección e integración del documento técnico del proyecto}
\textbf{Fecha:} 22/05/2026 \quad \textbf{Hora:} 19:00--23:30

\textbf{Objetivo:}
Unificar la información técnica del proyecto en un solo documento base y corregir errores de formato en LaTeX.

\textbf{Desarrollo:}
\begin{itemize}
\item Se revisaron documentos anteriores del proyecto para unir la información de sensores, acondicionamiento, alimentación, conectores y criterios de PCB.
\item Se depuró la estructura del documento principal, separando la información por bloques funcionales para facilitar su lectura.
\item Se corrigieron errores de compilación y de descarga del archivo, priorizando que el documento pudiera exportarse correctamente.
\item Se organizó una guía de continuación del proyecto, indicando qué partes faltaban por cerrar antes de pasar al esquemático definitivo.
\item Se reforzó que el diseño debe documentar no solo los cálculos, sino también las razones técnicas de cada decisión.
\end{itemize}

\textbf{Resultados:}
Se obtuvo una versión más ordenada del documento técnico del proyecto, lista para seguir siendo completada con los valores definitivos.

\textbf{Dificultades:}
La principal dificultad fue mantener toda la información en un solo archivo sin perder coherencia ni generar errores de compilación.

% -------------------------
\subsection{Actualización de tablas de sensores y criterios de acondicionamiento}
\textbf{Fecha:} 23/05/2026 \quad \textbf{Hora:} 10:00--15:00

\textbf{Objetivo:}
Actualizar las tablas del proyecto con datos más claros para el PT100 y el SM-50.

\textbf{Desarrollo:}
\begin{itemize}
\item Se actualizó la tabla de caracterización de sensores con los datos corregidos del SM-50 y del PT100.
\item Para el PT100 se reforzó el uso de una corriente de excitación de \SI{1}{mA}, debido a que permite convertir la resistencia del sensor en una tensión medible con bajo autocalentamiento.
\item Para el sensor de torque se mantuvo el criterio de puente resistivo con salida diferencial de baja amplitud, por lo que se requiere amplificación instrumental.
\item Se incluyeron recomendaciones prácticas de desacoplo, como capacitores de \SI{100}{nF} cercanos a los circuitos integrados y capacitores de mayor valor por zona de alimentación.
\item Se ajustó el lenguaje de las recomendaciones para que quedara más claro y defendible en una presentación oral.
\end{itemize}

\textbf{Resultados:}
Las tablas quedaron más alineadas con los sensores reales y con las necesidades de acondicionamiento de la PCB.

\textbf{Dificultades:}
Evitar que la tabla mezclara datos de cálculo preliminar con valores ya corregidos o confirmados por hoja técnica.

% =========================
\section{Undécima semana}

% -------------------------
\subsection{Memoria de cálculo del canal PT100 y del canal de torque}
\textbf{Fecha:} 27/05/2026 \quad \textbf{Hora:} 00:30--04:30

\textbf{Objetivo:}
Construir una memoria de cálculo para justificar la ganancia, sensibilidad y resolución de los dos canales de medición.

\textbf{Desarrollo:}
\begin{itemize}
\item Se desarrolló el modelo lineal del PT100 usando \(R_0=\SI{100}{\ohm}\) y una sensibilidad aproximada de \SI{0.385}{\ohm/\celsius}.
\item Se calculó que, con \SI{1}{mA} de excitación, el PT100 entrega aproximadamente \SI{100}{mV} a \SI{0}{\celsius} y \SI{138.5}{mV} a \SI{100}{\celsius}.
\item Se identificó que los \SI{100}{mV} iniciales corresponden a offset y que la información útil del rango de \SI{0}{\celsius} a \SI{100}{\celsius} está en la variación de \SI{38.5}{mV}.
\item Se calculó una ganancia corregida cercana a \(G\approx 86\) para usar el rango útil después de restar el offset.
\item Para torque se calculó que la sensibilidad del puente, con \SI{10}{V} de excitación y \SI{3}{mV/V}, produce una salida máxima cercana a \SI{30}{mV}.
\item Se verificó que una ganancia cercana a \(G\approx 100\) permite llevar la señal del SM-50 a un rango cómodo de adquisición sin saturar.
\end{itemize}

\textbf{Resultados:}
Se estableció una memoria de cálculo base para defender por qué el PT100 y el SM-50 requieren etapas de amplificación distintas.

\textbf{Dificultades:}
La dificultad principal fue separar correctamente el offset del PT100 de la señal útil de temperatura para no calcular una ganancia incorrecta.

% -------------------------
\subsection{Verificación de resolución y criterio de bits del sistema de adquisición}
\textbf{Fecha:} 27/05/2026 \quad \textbf{Hora:} 12:00--15:00

\textbf{Objetivo:}
Aclarar el papel del número de bits del sistema de adquisición en el proyecto.

\textbf{Desarrollo:}
\begin{itemize}
\item Se revisó que el número de bits no se estaba usando para seleccionar un ADC externo, ya que la digitalización final la realiza el sistema BoomBox.
\item Se interpretó el cálculo de bits como una verificación de resolución, no como una etapa de selección de convertidor.
\item Se reforzó que el diseño analógico debe entregar una señal limpia, escalada y compatible con el rango de entrada del sistema de adquisición.
\item Se aclaró que el objetivo es comprobar si la resolución disponible del sistema es suficiente para leer los cambios de torque y temperatura sin perder información importante.
\item Se relacionó esta verificación con el rango de salida del acondicionamiento, evitando usar ganancias excesivas solo por intentar llenar todo el rango de conversión.
\end{itemize}

\textbf{Resultados:}
Se aclaró que el cálculo de bits funciona como una comprobación de coherencia del diseño y no como una selección de hardware adicional.

\textbf{Dificultades:}
Diferenciar entre diseñar un ADC propio y acondicionar señales para un sistema de adquisición que ya posee su propia etapa de digitalización.

% -------------------------
\subsection{Selección preliminar de componentes de acondicionamiento}
\textbf{Fecha:} 27/05/2026 \quad \textbf{Hora:} 18:00--21:30

\textbf{Objetivo:}
Comparar opciones de amplificadores de instrumentación y operacionales para justificar la selección de componentes.

\textbf{Desarrollo:}
\begin{itemize}
\item Se compararon amplificadores de instrumentación adecuados para señales diferenciales pequeñas, dando prioridad a CMRR, offset, deriva, facilidad de ajuste de ganancia y disponibilidad.
\item Se revisó el uso del INA828 como opción principal por su precisión y por ser adecuado para señales de bajo nivel provenientes de sensores de puente y RTD.
\item Se analizaron alternativas de menor costo o más accesibles, como INA826 o AD620, para tener opciones defendibles si cambia la disponibilidad de componentes.
\item Se revisó el papel de amplificadores operacionales como LF353, entendiendo que no sustituyen directamente a un amplificador de instrumentación para la primera etapa de medida diferencial.
\item Se separó la comparación en dos familias: amplificadores de instrumentación para la extracción diferencial y operacionales para filtrado, buffer, offset o etapas auxiliares.
\end{itemize}

\textbf{Resultados:}
Se definió que la primera etapa crítica debe ser un amplificador de instrumentación, mientras que los operacionales quedan como apoyo para filtros, referencias o adaptación de salida.

\textbf{Dificultades:}
Evitar comparar directamente componentes que cumplen funciones distintas dentro de la cadena analógica.

% -------------------------
\subsection{Revisión de documentación imperix y naturaleza de la salida hacia BoomBox}
\textbf{Fecha:} 28/05/2026 \quad \textbf{Hora:} 13:00--15:00

\textbf{Objetivo:}
Aclarar cómo debe conectarse la PCB al ecosistema imperix BoomBox.

\textbf{Desarrollo:}
\begin{itemize}
\item Se revisó documentación relacionada con módulos imperix para entender la filosofía de entrada y salida analógica por conectores tipo RJ45.
\item Se corrigió la idea de que el enlace hacia la BoomBox fuese digital o por fibra en esta etapa; la salida de la PCB se mantiene como señal analógica acondicionada.
\item Se adoptó como criterio que la señal hacia la BoomBox debe ser diferencial, balanceada y transportada por cable STP/RJ45.
\item Se identificó que la BoomBox es la encargada de realizar la digitalización, por lo que la PCB debe concentrarse en protección, acondicionamiento, filtrado y adaptación de rango.
\item Se revisó que la salida debe mantenerse dentro de un margen seguro, evitando saturación del frente analógico de entrada.
\end{itemize}

\textbf{Resultados:}
Se corrigió la arquitectura de comunicación: sensor, PCB propia, salida analógica diferencial por STP/RJ45 y entrada analógica de la BoomBox.

\textbf{Dificultades:}
Interpretar correctamente el uso físico del RJ45, ya que no se trata de Ethernet sino de un conector multipar para señales analógicas y alimentación.

% =========================
\section{Duodécima semana}

% -------------------------
\subsection{Revisión de alimentación para pruebas y protoboard}
\textbf{Fecha:} 31/05/2026 \quad \textbf{Hora:} 16:00--18:00

\textbf{Objetivo:}
Definir la forma más recomendable de alimentar el circuito durante pruebas y en la versión final.

\textbf{Desarrollo:}
\begin{itemize}
\item Se revisó que el circuito requiere rieles de alimentación compatibles con etapas analógicas, principalmente \(\pm\SI{15}{V}\) para permitir margen de operación.
\item Se diferenció la alimentación de pruebas en protoboard de la alimentación final de la PCB.
\item Se reforzó que la protoboard puede servir para validar funciones básicas, pero no es una prueba definitiva de ruido debido a capacitancias parásitas, falta de plano de tierra y conexiones largas.
\item Se revisó que la alimentación final no debe provenir de los sensores, sino del sistema de adquisición o de la interfaz definida para la tarjeta.
\item Se consideró la necesidad de desacoplo local por circuito integrado y distribución ordenada de los rieles dentro de la PCB.
\end{itemize}

\textbf{Resultados:}
Se estableció que las pruebas pueden hacerse con fuente de laboratorio, pero el diseño final debe considerar alimentación estable, desacoplada y correctamente distribuida.

\textbf{Dificultades:}
Separar lo que es válido para una prueba rápida en protoboard de lo que debe cumplir una PCB final cercana a un motor e inversor.

% -------------------------
\subsection{Comparación técnica de etapas diferenciales e instrumentales}
\textbf{Fecha:} 31/05/2026 \quad \textbf{Hora:} 20:00--23:30

\textbf{Objetivo:}
Comparar técnicamente amplificadores diferenciales y de instrumentación para cerrar la arquitectura del acondicionamiento.

\textbf{Desarrollo:}
\begin{itemize}
\item Se revisaron fuentes técnicas sobre acondicionamiento de puentes Wheatstone y medición de señales diferenciales de baja amplitud.
\item Se comparó el uso de un amplificador diferencial con resistencias externas frente a un amplificador de instrumentación integrado.
\item Se identificó que el amplificador diferencial exige resistencias muy bien apareadas para mantener buen rechazo de modo común.
\item Se concluyó que el amplificador de instrumentación es más robusto para la primera etapa porque ofrece alta impedancia de entrada y alto CMRR.
\item Se dejó el operacional diferencial como opción para etapas posteriores, como adaptación de salida, offset, filtrado activo o conversión diferencial si se requiere.
\end{itemize}

\textbf{Resultados:}
Se fortaleció la decisión de usar un INA como primera etapa para torque y temperatura, dejando los operacionales como bloques auxiliares.

\textbf{Dificultades:}
Defender por qué dos circuitos que parecen similares en esquema no se comportan igual frente a modo común, tolerancias y ruido.

% -------------------------
\subsection{Definición de conectores del SM-50 y salidas RJ45}
\textbf{Fecha:} 01/06/2026 \quad \textbf{Hora:} 02:00--05:00

\textbf{Objetivo:}
Definir conectores prácticos para el sensor SM-50 y para la interfaz de la PCB con la BoomBox.

\textbf{Desarrollo:}
\begin{itemize}
\item Se revisó el conector físico del sensor SM-50 y la necesidad de adaptar sus líneas hacia una entrada robusta en PCB.
\item Se propuso usar una bornera push-in de 5 pines para recibir +EXC, -EXC, +SIG, -SIG y SHIELD.
\item Se planteó un arnés adaptador desde el conector del sensor hacia ferrules y bornera, evitando modificar el cable original del sensor.
\item Se definió que la salida hacia la BoomBox debe usar conectores RJ45 blindados, uno para torque y otro para temperatura.
\item Se estableció que la pantalla del RJ45 debe tratarse como blindaje o chasis, no como una conexión directa e indiscriminada al GND analógico.
\end{itemize}

\textbf{Resultados:}
Se obtuvo una propuesta práctica de conectores: entrada SM-50 por bornera de 5 pines y salida hacia BoomBox por RJ45 blindado.

\textbf{Dificultades:}
Elegir conectores que fueran fáciles de probar en laboratorio pero suficientemente robustos para una PCB real.

% -------------------------
\subsection{Organización de la PCB por ciudades funcionales}
\textbf{Fecha:} 01/06/2026 \quad \textbf{Hora:} 20:00--23:30

\textbf{Objetivo:}
Organizar físicamente la PCB por bloques o ciudades para reducir ruido y facilitar el diseño.

\textbf{Desarrollo:}
\begin{itemize}
\item Se definió una organización por ciudades para separar físicamente las zonas sensibles, la alimentación y las salidas.
\item La ciudad 1 corresponde a la entrada del SM-50, excitación del puente y filtro inicial.
\item La ciudad 2 corresponde al acondicionamiento de torque mediante INA, ajuste de ganancia y filtrado.
\item La ciudad 3 corresponde a la entrada del PT100 mediante bornera de 4 hilos, fuente de corriente y líneas de sense.
\item La ciudad 4 corresponde al acondicionamiento de temperatura, resta de offset, ganancia y filtrado.
\item La ciudad 5 corresponde a la alimentación desde la BoomBox, protección y desacoplo.
\item La ciudad 6 corresponde a las salidas RJ45 hacia la BoomBox para torque y temperatura.
\end{itemize}

\textbf{Resultados:}
Se obtuvo una distribución funcional que permite diseñar el esquemático y el layout con mejor control del ruido.

\textbf{Dificultades:}
Evitar mezclar físicamente señales de milivoltios con alimentación, conectores de salida y posibles retornos ruidosos.

% -------------------------
\subsection{Ajuste de alimentación final desde B-Box y referencia de \SI{3.3}{V}}
\textbf{Fecha:} 02/06/2026 \quad \textbf{Hora:} 13:00--18:00

\textbf{Objetivo:}
Cerrar la estrategia de alimentación de la PCB y justificar el uso de rieles internos auxiliares.

\textbf{Desarrollo:}
\begin{itemize}
\item Se revisó que la PCB no usará una fuente externa como alimentación principal, sino los rieles disponibles desde la B-Box por medio de los puertos analógicos RJ45.
\item Se definió que los rieles principales de entrada serán \( +\SI{15}{V} \), \( -\SI{15}{V} \) y GND, distribuidos por los conectores RJ45.
\item Se analizó que algunas etapas analógicas pueden trabajar directamente con rieles cercanos a \(\pm\SI{15}{V}\), mientras que referencias internas o bloques auxiliares pueden requerir \SI{3.3}{V}.
\item Se revisó el concepto de LDO como regulador lineal de baja caída para obtener una referencia o alimentación auxiliar limpia de \SI{3.3}{V}.
\item Se consideró la corriente disponible aproximada de \SI{100}{mA}, por lo que el consumo de referencias, INAs, operacionales y excitaciones debe mantenerse bajo control.
\item Se reforzó que la salida final debe aprovechar el rango permitido por la BoomBox sin sacrificar margen frente a saturación.
\end{itemize}

\textbf{Resultados:}
Se definió una arquitectura de alimentación más coherente: entrada \(\pm\SI{15}{V}\) desde B-Box, regulación local cuando haga falta y desacoplo por bloques.

\textbf{Dificultades:}
Distinguir entre alimentación principal, referencia de medición, excitación de sensores y niveles internos auxiliares.

% -------------------------
\subsection{Diseño conceptual de excitación para PT100 de tres hilos}
\textbf{Fecha:} 03/06/2026 \quad \textbf{Hora:} 15:30--20:30

\textbf{Objetivo:}
Definir una forma más segura y precisa de excitar el PT100 usando \SI{1}{mA}.

\textbf{Desarrollo:}
\begin{itemize}
\item Se revisó el uso de PT100 de tres hilos para compensar la resistencia de los cables cuando el sensor se encuentra separado de la PCB.
\item Se comparó la opción de una fuente de corriente hecha con amplificador operacional y transistor frente al uso de un circuito integrado de corriente de precisión.
\item Se identificó que una solución integrada, como REF200 o una fuente de corriente dedicada, reduce el riesgo de errores de implementación frente a una etapa discreta con LM358 y transistor.
\item Se mantuvo el criterio de \SI{1}{mA} como corriente de excitación para obtener una señal medible y reducir el autocalentamiento.
\item Se revisó que la tensión del PT100 debe medirse de forma diferencial y que el offset de \SI{100}{mV} a \SI{0}{\celsius} debe tratarse con cuidado antes de aplicar alta ganancia.
\item Se reforzó que el ajuste de calibración debe realizarse con trimmer cuando sea necesario, no con un potenciómetro de panel como elemento principal del diseño.
\end{itemize}

\textbf{Resultados:}
Se orientó el canal PT100 hacia una fuente de corriente más precisa y menos propensa a fallos, manteniendo la medición diferencial y la compensación de cableado.

\textbf{Dificultades:}
Elegir una arquitectura que sea precisa, fácil de defender y al mismo tiempo tenga baja probabilidad de fallar en la implementación física.

% =========================
\section{Decimotercera semana}

% -------------------------
\subsection{Revisión de documentación tipo handbook para esquemáticos de PCB}
\textbf{Fecha:} 07/06/2026 \quad \textbf{Hora:} 17:30--19:00

\textbf{Objetivo:}
Buscar referencias de diseño esquemático para orientar la presentación del circuito final.

\textbf{Desarrollo:}
\begin{itemize}
\item Se revisó la necesidad de usar ejemplos tipo handbook o notas de aplicación para que el esquemático no dependa únicamente de criterio propio.
\item Se identificaron fuentes útiles de fabricantes y notas de aplicación para puentes resistivos, RTD, acondicionamiento de señal y diseño de PCB.
\item Se relacionaron esas fuentes con los bloques del proyecto: excitación, protección, amplificación instrumental, filtrado, referencias y salida diferencial.
\item Se priorizó usar esquemas de aplicación como guía conceptual, adaptándolos a los sensores reales y al ecosistema BoomBox.
\end{itemize}

\textbf{Resultados:}
Se reforzó la estrategia de respaldar el diseño con circuitos de referencia y no solo con cálculos aislados.

\textbf{Dificultades:}
Encontrar ejemplos que sean parecidos al proyecto sin copiar una aplicación que no coincida con los rangos, sensores o alimentación disponibles.

% -------------------------
\subsection{Cierre de arquitectura final para PT100 y SM-50}
\textbf{Fecha:} 07/06/2026 \quad \textbf{Hora:} 19:30--22:30

\textbf{Objetivo:}
Actualizar la arquitectura final del informe considerando las últimas decisiones del proyecto.

\textbf{Desarrollo:}
\begin{itemize}
\item Se consolidó la arquitectura final separando claramente el canal PT100 y el canal SM-50.
\item Para el canal PT100 se decidió priorizar una fuente de corriente de precisión de \SI{1}{mA}, medición diferencial, resta de offset, ganancia moderada y filtrado pasabajo.
\item Se revisó la posibilidad de eliminar una fuente de corriente discreta con LM358 y transistor cuando una fuente integrada ofrezca mayor precisión y menor riesgo de error práctico.
\item Para el canal SM-50 se mantuvo la excitación del puente con referencia estable de \SI{10}{V}, lectura diferencial del puente, amplificación instrumental y filtrado.
\item Se corrigió la idea de usar potenciómetro, dejando claro que el ajuste físico adecuado en la PCB debe ser mediante trimmer multivuelta si se requiere calibración fina.
\item Se reforzó que el informe final debe distinguir entre alimentación de la PCB, excitación de sensores, referencia de offset, ganancia del INA y salida hacia la BoomBox.
\item Se agregó como criterio final que la PCB debe ser robusta antes que excesivamente compleja: menos bloques discretos cuando un integrado confiable simplifica el diseño.
\end{itemize}

\textbf{Resultados:}
La arquitectura final quedó más clara: PT100 con excitación precisa de \SI{1}{mA} y acondicionamiento diferencial; SM-50 con excitación estable de \SI{10}{V} y amplificación instrumental.

\textbf{Dificultades:}
Ajustar el informe para que no quedaran mezcladas decisiones preliminares con decisiones finales ya corregidas.% =========================
\section{Decimocuarta semana}

% -------------------------
\subsection{Depuración de arquitectura y fuente de corriente del PT100}
\textbf{Fecha:} 08/06/2026 \quad \textbf{Hora:} 16:00--20:30

\textbf{Objetivo:}
Revisar la arquitectura del canal de temperatura y reducir riesgos de implementación en la fuente de corriente del PT100.

\textbf{Desarrollo:}
\begin{itemize}
\item Se revisó nuevamente la necesidad de excitar el PT100 con una corriente estable para convertir la variación de resistencia en una tensión proporcional a la temperatura.
\item Se comparó una fuente de corriente discreta basada en operacional, transistor y resistencia de ajuste frente a una alternativa integrada o una referencia de precisión con resistencia de programación.
\item Se identificó que una solución con menos bloques discretos facilita la explicación del diseño, reduce puntos de falla y mejora la repetibilidad en PCB.
\item Se mantuvo la prioridad de medir el PT100 en forma diferencial para aprovechar el rechazo de modo común y reducir errores por cableado.
\item Se reforzó que el canal de temperatura debe separar claramente tres conceptos: excitación del sensor, resta de offset y amplificación hacia el rango permitido por la adquisición.
\end{itemize}

\textbf{Resultados:}
Se consolidó la idea de usar una fuente de corriente precisa para el PT100 y mantener la medición diferencial como base del canal de temperatura.

\textbf{Dificultades:}
Evitar que el esquemático creciera innecesariamente con bloques discretos difíciles de ajustar y defender durante la exposición.

% -------------------------
\subsection{Revisión de referencias técnicas para RTD, puente Wheatstone y acondicionamiento}
\textbf{Fecha:} 10/06/2026 \quad \textbf{Hora:} 19:00--22:30

\textbf{Objetivo:}
Respaldar las decisiones del diseño con documentación técnica sobre sensores resistivos, RTD y puentes Wheatstone.

\textbf{Desarrollo:}
\begin{itemize}
\item Se revisaron notas de aplicación y material técnico sobre acondicionamiento de sensores resistivos y medición de señales diferenciales de baja amplitud.
\item Para el SM-50 se reforzó que el sensor se comporta como un puente Wheatstone de galgas extensométricas con salida de pocos milivoltios por voltio de excitación.
\item Se relacionó la necesidad del amplificador de instrumentación con la alta impedancia de entrada, el alto rechazo de modo común y la baja amplitud de la señal diferencial.
\item Para el PT100 se revisó el modelo resistivo del RTD y la conveniencia de una fuente de corriente para obtener una relación directa entre resistencia y tensión.
\item Se identificaron como temas clave para justificar el diseño: sensibilidad, offset, ganancia, filtrado, rango dinámico, protección y ruido.
\end{itemize}

\textbf{Resultados:}
Se obtuvo una base teórica más sólida para explicar por qué ambos canales requieren acondicionamiento analógico antes de la entrada diferencial de la BoomBox.

\textbf{Dificultades:}
Integrar referencias de fabricantes y teoría de sensores sin que el informe se volviera una copia de notas de aplicación.

% -------------------------
\subsection{Ajuste de alcance del proyecto y bloques mínimos de la PCB}
\textbf{Fecha:} 14/06/2026 \quad \textbf{Hora:} 15:00--19:30

\textbf{Objetivo:}
Ajustar el alcance práctico de la PCB para que el diseño fuera fabricable, explicable y compatible con el tiempo disponible.

\textbf{Desarrollo:}
\begin{itemize}
\item Se revisó la diferencia entre un diseño ampliado con varios canales y una versión mínima funcional para validar el principio de medición.
\item Se priorizó trabajar con un canal de PT100 y un canal de SM-50, manteniendo la posibilidad de replicar bloques si se requiere escalar el diseño.
\item Se separaron los bloques indispensables: entrada de sensor, protección, excitación, amplificación instrumental, filtrado de salida, alimentación y conectores hacia la BoomBox.
\item Se revisó que la PCB no debe digitalizar la señal, sino acondicionarla para que la adquisición se realice externamente.
\item Se identificó que reducir el alcance ayuda a mejorar el ruteo, el orden físico y la revisión de errores en KiCad.
\end{itemize}

\textbf{Resultados:}
El proyecto quedó enfocado en una PCB analógica de acondicionamiento para un PT100 y un SM-50, con bloques replicables y mejor control de complejidad.

\textbf{Dificultades:}
Recortar el alcance sin perder coherencia técnica ni dejar fuera protecciones necesarias para un entorno con motor e inversor.

% =========================
\section{Decimoquinta semana}

% -------------------------
\subsection{Simulación en LTspice y revisión de memoria de cálculo}
\textbf{Fecha:} 17/06/2026 \quad \textbf{Hora:} 14:00--20:00

\textbf{Objetivo:}
Verificar mediante simulación y cálculo los rangos esperados de los canales PT100 y SM-50.

\textbf{Desarrollo:}
\begin{itemize}
\item Se trabajó con esquemáticos de simulación para el canal PT100 y el canal SM-50 con el fin de comprobar niveles de señal antes de trasladarlos a KiCad.
\item En el canal PT100 se revisó la relación entre corriente de excitación, resistencia del sensor, tensión diferencial y ganancia necesaria.
\item Se comparó mantener \SI{1}{mA} frente a reducir la excitación del PT100 a aproximadamente \SI{0.5}{mA}, buscando menor autocalentamiento y mejor margen de consumo.
\item En el canal SM-50 se revisó la excitación del puente, la salida diferencial esperada y la ganancia del amplificador de instrumentación.
\item Se ajustaron valores de resistencias de ganancia para mantener la salida dentro del rango seguro de la adquisición.
\item Se revisó la memoria de cálculo para corregir decisiones preliminares y dejar coherencia entre cálculo, simulación y esquemático.
\end{itemize}

\textbf{Resultados:}
Se validó que los dos canales pueden entregar señales acondicionadas dentro del margen requerido, siempre que se controle la ganancia y el filtrado de salida.

\textbf{Dificultades:}
Mantener consistentes los valores entre LTspice, memoria de cálculo, lista de materiales y esquemático de KiCad.

% -------------------------
\subsection{Selección de amplificadores y reducción de costo del diseño}
\textbf{Fecha:} 18/06/2026 \quad \textbf{Hora:} 18:30--22:30

\textbf{Objetivo:}
Comparar opciones de amplificadores para lograr una solución precisa sin elevar innecesariamente el costo de la PCB.

\textbf{Desarrollo:}
\begin{itemize}
\item Se comparó el uso de INA828 frente a INA826 para los canales de acondicionamiento.
\item Se identificó que el INA826 representa una alternativa más económica para mantener una topología de instrumentación adecuada.
\item Se revisó el uso del OPA2188 como operacional de precisión para etapas auxiliares, principalmente donde el offset y la deriva pueden afectar la medición.
\item Se discutió que el PT100 no exige una arquitectura excesivamente compleja si la señal se acondiciona de forma limpia y se calibra correctamente.
\item Se revisó la necesidad de no sacrificar rechazo de modo común ni estabilidad por ahorrar componentes críticos.
\end{itemize}

\textbf{Resultados:}
Se orientó el diseño hacia el uso de amplificadores de instrumentación INA826 y operacionales de precisión para las funciones auxiliares necesarias.

\textbf{Dificultades:}
Balancear costo, disponibilidad, precisión y facilidad de justificación técnica.

% -------------------------
\subsection{Investigación de alimentación, LDO y límite de corriente de la B-Box}
\textbf{Fecha:} 19/06/2026 \quad \textbf{Hora:} 13:00--21:00

\textbf{Objetivo:}
Definir una alimentación realista para la PCB considerando que la B-Box entrega rieles limitados en corriente.

\textbf{Desarrollo:}
\begin{itemize}
\item Se revisó que la alimentación disponible desde la B-Box es de \(\pm\SI{15}{V}\) con una corriente limitada, por lo que el consumo total de la PCB debe mantenerse bajo control.
\item Se investigaron reguladores lineales tipo LDO en encapsulado SMD para generar una tensión limpia a partir del riel de entrada.
\item Se compararon opciones como reguladores de la familia 78L05/78M05 y reguladores de mayor calidad, considerando disponibilidad, ruido, corriente y facilidad de montaje.
\item Se definió que la excitación del SM-50 debe ser estable y de bajo ruido, debido a que cualquier variación de excitación afecta directamente la salida del puente.
\item Se revisó la opción de obtener \SI{10}{V} para el SM-50 desde el riel de \SI{15}{V} mediante regulación local, manteniendo la excitación separada de la señal de salida.
\item Se reforzó la necesidad de colocar capacitores de desacoplo cerca de los reguladores, referencias y amplificadores.
\end{itemize}

\textbf{Resultados:}
Se estableció que la alimentación debe partir de los rieles de la B-Box, con regulación local y desacoplo suficiente para referencias, excitación del SM-50 y fuente de corriente del PT100.

\textbf{Dificultades:}
Elegir un regulador disponible y defendible sin exceder la corriente disponible ni introducir ruido en las señales de milivoltios.

% -------------------------
\subsection{Revisión de EMI, filtros y necesidad de protección en la PCB}
\textbf{Fecha:} 20/06/2026 \quad \textbf{Hora:} 16:00--20:00

\textbf{Objetivo:}
Definir medidas básicas contra ruido e interferencias para una PCB cercana a un sistema de potencia.

\textbf{Desarrollo:}
\begin{itemize}
\item Se revisó que el entorno del motor y del inversor puede introducir interferencia electromagnética por acoplamiento capacitivo, inductivo y conducido.
\item Se discutió el uso de filtros RC en entradas y salidas para limitar ruido de alta frecuencia sin afectar la banda útil de torque y temperatura.
\item Se reforzó el uso de señales diferenciales y cableado blindado para reducir sensibilidad al ruido externo.
\item Se revisó que el blindaje de cables y conectores debe tratarse con cuidado para no contaminar el GND analógico sensible.
\item Se consideró la necesidad de TVS, fusibles rearmables y desacoplos como parte del diseño robusto, no como componentes opcionales decorativos.
\end{itemize}

\textbf{Resultados:}
Se definió que la PCB debe incorporar filtrado, protección y una organización física que separe sensores, acondicionamiento, alimentación y salidas.

\textbf{Dificultades:}
Evitar llenar el esquemático de filtros innecesarios, pero sin dejar expuestas entradas y alimentaciones sensibles.

% -------------------------
\subsection{Integración del esquemático en KiCad y corrección de errores ERC}
\textbf{Fecha:} 21/06/2026 \quad \textbf{Hora:} 10:00--22:00

\textbf{Objetivo:}
Avanzar la integración del esquemático final en KiCad y resolver errores eléctricos antes de pasar al layout.

\textbf{Desarrollo:}
\begin{itemize}
\item Se revisaron símbolos, footprints y etiquetas de red para que las etapas quedaran conectadas correctamente.
\item Se corrigieron problemas relacionados con pines de alimentación, referencias y uso de \texttt{PWR\_FLAG} en KiCad.
\item Se definieron etiquetas claras para unir bloques de salida, resistencias limitadoras, trimmers, resistencias de aislamiento y salidas diferenciales.
\item Se revisó la asignación de footprints para componentes SMD principales, incluyendo amplificadores, referencia, reguladores, capacitores y resistencias.
\item Se depuraron componentes no utilizados para reducir ruido visual en el esquemático y evitar confusiones en la lista de materiales.
\end{itemize}

\textbf{Resultados:}
El esquemático quedó más ordenado y preparado para pasar al diseño físico, con conexiones más claras entre bloques funcionales.

\textbf{Dificultades:}
Resolver errores de KiCad sin esconder problemas reales de alimentación o conexión entre bloques.

% =========================
\section{Decimosexta semana}

% -------------------------
\subsection{Selección e importación de protecciones TVS y fusibles rearmables}
\textbf{Fecha:} 22/06/2026 \quad \textbf{Hora:} 09:00--18:00

\textbf{Objetivo:}
Agregar protecciones eléctricas al diseño y asegurar que los símbolos y huellas fueran compatibles con fabricación.

\textbf{Desarrollo:}
\begin{itemize}
\item Se revisó el uso de TVS para proteger líneas expuestas frente a eventos transitorios y descargas electrostáticas.
\item Se incorporó la protección TVS de \SI{12}{V} para la zona de entrada de sensores y se revisó su conexión respecto a la señal y tierra de referencia.
\item Se revisaron TVS/ESD de \SI{5.5}{V} para líneas sensibles hacia la interfaz de adquisición, especialmente donde se requiere proteger entradas de bajo nivel.
\item Se incluyeron fusibles rearmables PPTC en las entradas de alimentación para limitar fallas por sobrecorriente o errores de conexión.
\item Se importaron y verificaron símbolos y footprints, corrigiendo pines y encapsulados cuando no coincidían con el componente real.
\end{itemize}

\textbf{Resultados:}
El diseño incorporó protecciones más robustas mediante TVS y PPTC, manteniendo una organización por bloques para facilitar revisión y ruteo.

\textbf{Dificultades:}
Diferenciar qué protección corresponde a alimentación, a entradas de sensores y a líneas de salida hacia la BoomBox.

% -------------------------
\subsection{Conectores, RJ45 blindados y organización de ciudades de la PCB}
\textbf{Fecha:} 22/06/2026 \quad \textbf{Hora:} 19:00--23:30

\textbf{Objetivo:}
Definir el acomodo físico de conectores y bloques funcionales antes de comenzar el ruteo principal.

\textbf{Desarrollo:}
\begin{itemize}
\item Se mantuvo el uso de conectores RJ45 blindados para la conexión con la BoomBox y cable FTP CAT5.
\item Se revisó la asignación de pines de la interfaz RJ45 considerando rieles de \(\pm\SI{15}{V}\), GND y entrada diferencial.
\item Se separó el tratamiento del blindaje respecto al GND analógico para evitar que corrientes de pantalla circulen por rutas sensibles.
\item Se ordenó el layout por ciudades: entradas de sensores, canal SM-50, canal PT100, alimentación/protecciones, salidas RJ45 y zona de documentación.
\item Se definió que los conectores de sensores deben ubicarse cerca de sus protecciones, mientras que los conectores hacia BoomBox deben quedar del lado de salida.
\end{itemize}

\textbf{Resultados:}
La PCB quedó organizada con una lógica física más clara, reduciendo cruces de señales y separando zonas ruidosas de zonas sensibles.

\textbf{Dificultades:}
Ubicar conectores grandes, borneras, fusibles, TVS y circuitos integrados sin romper la separación entre señales de entrada, alimentación y salida.

% -------------------------
\subsection{Ruteo preliminar y definición de anchos de pista}
\textbf{Fecha:} 24/06/2026 \quad \textbf{Hora:} 09:00--21:00

\textbf{Objetivo:}
Iniciar el ruteo de la PCB usando reglas prácticas para señales sensibles, alimentación y conexiones de salida.

\textbf{Desarrollo:}
\begin{itemize}
\item Se revisó el ruteo por ciudades para mantener señales de milivoltios lejos de reguladores, conectores de alimentación y retornos con corriente.
\item Se asignaron anchos de pista distintos para señales analógicas, alimentación, excitación de sensores y salidas hacia conectores.
\item Se evaluó cuándo usar una segunda capa para evitar cruces y mantener rutas cortas, especialmente en zonas congestionadas cerca de conectores.
\item Se revisó la necesidad de colocar capacitores de desacoplo cerca de cada circuito integrado y de mantener los filtros cerca de los pines que protegen.
\item Se corrigieron problemas de elementos bloqueados en KiCad y se mantuvieron fijos el marco y los agujeros de montaje.
\end{itemize}

\textbf{Resultados:}
Se obtuvo un ruteo preliminar más ordenado, con criterios diferenciados para pistas de señal, alimentación y protección.

\textbf{Dificultades:}
El mayor problema fue la congestión cerca de conectores y componentes de protección, donde no siempre era posible usar el ancho de pista ideal.

% -------------------------
\subsection{Revisión de footprints mecánicos y dimensiones de borneras}
\textbf{Fecha:} 25/06/2026 \quad \textbf{Hora:} 10:00--13:00

\textbf{Objetivo:}
Verificar que los componentes mecánicos principales tuvieran espacio suficiente en la PCB.

\textbf{Desarrollo:}
\begin{itemize}
\item Se revisaron las dimensiones físicas de las borneras utilizadas para entrada de sensores.
\item Se buscó definir el rectángulo de ocupación de cada bornera para dejar espacio alrededor y evitar interferencia con otros componentes.
\item Se consideró que las borneras deben quedar accesibles para cableado, ferrules y manipulación durante pruebas de laboratorio.
\item Se verificó que los conectores grandes condicionan el borde de la PCB y el orden de las ciudades cercanas.
\item Se mantuvo el criterio de dejar los componentes de protección cerca de los conectores de entrada.
\end{itemize}

\textbf{Resultados:}
Se mejoró la planificación mecánica del layout, especialmente en la ubicación de borneras, conectores y protecciones de entrada.

\textbf{Dificultades:}
Ajustar componentes grandes en una placa compacta sin perder claridad de ruteo ni accesibilidad para pruebas.

% -------------------------
\subsection{Preparación del póster electrónico y síntesis de avances}
\textbf{Fecha:} 25/06/2026 \quad \textbf{Hora:} 15:00--18:30

\textbf{Objetivo:}
Preparar la información del proyecto para la exposición tipo póster electrónico de IE0499.

\textbf{Desarrollo:}
\begin{itemize}
\item Se revisó la rúbrica de exposición para priorizar avances, dificultades y aprendizaje incorporado durante el desarrollo del proyecto.
\item Se definió que el póster debe ser horizontal, legible en pantalla y con máximo dos transparencias.
\item Se organizaron los contenidos principales: problema, objetivo, arquitectura de la PCB, avances técnicos, dificultades y conclusiones.
\item Se seleccionaron como avances principales la definición de los canales PT100 y SM-50, la integración en KiCad, la selección de protecciones y el inicio del ruteo.
\item Se seleccionaron como dificultades principales el control de ruido, la limitación de corriente de la B-Box, la importación de footprints y la congestión del layout.
\end{itemize}

\textbf{Resultados:}
Se obtuvo una estructura de póster más enfocada en explicar decisiones técnicas y dificultades reales, no solo en mostrar el circuito terminado.

\textbf{Dificultades:}
Sintetizar varias semanas de trabajo en pocas ideas claras que puedan defenderse durante preguntas de otras áreas.



% =========================
\section{Decimoséptima semana}

% -------------------------
\subsection{Actualización final de valores del canal SM-50 y PT100}
\textbf{Fecha:} 26/06/2026 \quad \textbf{Hora:} 18:00--23:30

\textbf{Objetivo:}
Actualizar los cálculos del proyecto con los valores finales realmente usados en la PCB y separar las decisiones vigentes de las memorias antiguas.

\textbf{Desarrollo:}
\begin{itemize}
\item Se revisaron los cálculos finales del canal de torque SM-50, manteniendo la excitación del puente en \(\SI{10}{V}\) y usando la sensibilidad nominal de \(\SI{3}{mV/V}\), por lo que la señal diferencial máxima esperada del puente es aproximadamente \(\SI{30}{mV}\).
\item Se actualizó la resistencia de ganancia del amplificador de instrumentación del canal SM-50 a \(R_G=\SI{160}{\ohm}\), con el fin de obtener una mayor ganancia y mejorar el aprovechamiento del rango de entrada de la BoomBox.
\item Se eliminó del diseño la resistencia en serie de \(\SI{180}{\ohm}\) que estaba considerada en versiones anteriores, ya que reducía el rango útil y no correspondía con la versión actual del acondicionamiento.
\item Se corrigió el canal de temperatura indicando que la excitación final del PT100 se toma como \(I_{exc}=\SI{0.5}{mA}\), no \(\SI{1}{mA}\), con el propósito de reducir autocalentamiento y mantener bajo el consumo.
\item Se mantuvo el criterio del PT100 de tres hilos, usando medición diferencial, compensación de cableado, referencia de offset y ajuste de nivel antes de entregar la señal a la BoomBox.
\item Se dejó como valor vigente para el canal PT100 una resistencia de ganancia del INA de \(R_G=\SI{1}{k\ohm}\), separando este dato de valores preliminares usados en cálculos anteriores.
\item Se verificó que ambos canales requieren ajuste de nivel/salida: el canal SM-50 para desplazar y escalar la salida hacia el rango útil de la BoomBox, y el canal PT100 para tratar el offset propio de la resistencia base de \(\SI{100}{\ohm}\).
\end{itemize}

\textbf{Resultados:}
Se actualizaron los valores finales del diseño: SM-50 con \(R_G=\SI{160}{\ohm}\), sin resistencia serie de \(\SI{180}{\ohm}\), salida aproximada de \(\SIrange{0}{10}{V}\); y PT100 con \(I_{exc}=\SI{0.5}{mA}\) y \(R_G=\SI{1}{k\ohm}\).

\textbf{Dificultades:}
La principal dificultad fue no mezclar valores de memorias antiguas con los datos finales del circuito, especialmente en la corriente de excitación del PT100 y en la ganancia del SM-50.

% -------------------------
\subsection{Validación de interfaz BoomBox, 8P8C y cable FTP Cat5}
\textbf{Fecha:} 26/06/2026 \quad \textbf{Hora:} 23:30--02:00

\textbf{Objetivo:}
Confirmar la conexión física y eléctrica entre la PCB y la BoomBox usando conectores 8P8C/RJ45 y cable FTP Cat5.

\textbf{Desarrollo:}
\begin{itemize}
\item Se revisó el manual de la BoomBox para confirmar que la entrada analógica puede configurarse en modo de alta impedancia diferencial de \(\SI{3}{k\ohm}\).
\item Se verificó que la BoomBox entrega alimentación para sensores externos mediante el conector analógico, con rieles de \(\pm\SI{15}{V}\) y corriente máxima de \(\SI{100}{mA}\).
\item Se actualizó la documentación para indicar que la conexión física entre la PCB y la BoomBox se realiza mediante cable FTP Cat5 y conector modular 8P8C blindado.
\item Se mantuvo la asignación de pines del conector: pines 1 y 2 para \(+\SI{15}{V}\), pines 3 y 6 para \(\SI{0}{V}\), pines 4 y 5 para señal analógica diferencial, y pines 7 y 8 para \(-\SI{15}{V}\).
\item Se aclaró que el conector 8P8C no se utiliza como Ethernet, sino como interfaz física multipar para alimentación y señal analógica.
\item Se definió que el par de señal debe mantenerse como par trenzado y que el blindaje del conector debe tratarse como pantalla, evitando mezclarlo directamente con el retorno analógico sensible.
\end{itemize}

\textbf{Resultados:}
Se corrigió la representación de la interfaz hacia BoomBox: entrada de alta impedancia diferencial, alimentación limitada a \(\pm\SI{15}{V}\) y \(\SI{100}{mA}\), transporte por FTP Cat5/8P8C y separación clara entre señal, alimentación y blindaje.

\textbf{Dificultades:}
Interpretar correctamente el conector RJ45, ya que físicamente se parece a una red Ethernet, pero en el proyecto funciona como interfaz analógica y de alimentación.

% -------------------------
\subsection{Separación de memorias antiguas y trazabilidad del proceso}
\textbf{Fecha:} 27/06/2026 \quad \textbf{Hora:} 00:00--03:00

\textbf{Objetivo:}
Organizar los documentos previos del proyecto para conservarlos como evidencia del proceso sin usar sus valores obsoletos como cálculos finales.

\textbf{Desarrollo:}
\begin{itemize}
\item Se revisaron memorias antiguas, avances preliminares y documentos de arquitectura generados durante el desarrollo del proyecto.
\item Se clasificó como información vigente únicamente lo relacionado con el pinout 8P8C/RJ45, el uso de cable FTP/STP, la separación de puertos por canal y la justificación conceptual de los sensores resistivos.
\item Se marcó como obsoleta la información asociada a ADC de \SI{12}{bits}, referencia de \SI{3.3}{V}, PT100 con \SI{1}{mA} como valor final, ganancias preliminares y BOM anteriores.
\item Se conservó la información teórica útil: modelo del PT100, conexión de tres hilos, fuente de corriente para RTD, puente Wheatstone completo, sensibilidad del SM-50 y necesidad de amplificación instrumental.
\item Se redactó una nota de trazabilidad indicando que los documentos anteriores explican la evolución del diseño, pero que la memoria final usa los valores actuales.
\end{itemize}

\textbf{Resultados:}
Se evitó que la memoria final mezclara cálculos viejos con el diseño actual. Las memorias antiguas quedaron como respaldo histórico del proceso de aprendizaje y toma de decisiones.

\textbf{Dificultades:}
Distinguir qué partes de documentos anteriores seguían siendo útiles como referencia conceptual y cuáles debían descartarse para no generar contradicciones técnicas.

% -------------------------
\subsection{Actualización de memoria de cálculo, referencias y diagrama final}
\textbf{Fecha:} 27/06/2026 \quad \textbf{Hora:} 03:00--06:00

\textbf{Objetivo:}
Actualizar la memoria de cálculo del proyecto con los valores actuales, referencias completas y diagramas finales.

\textbf{Desarrollo:}
\begin{itemize}
\item Se actualizó la memoria de cálculo del proyecto separando los cálculos finales de la trazabilidad de documentos previos.
\item Se incorporaron los valores actuales del canal PT100: \(I_{exc}=\SI{0.5}{mA}\), \(R_G=\SI{1}{k\ohm}\), conexión de tres hilos, referencia de offset proveniente de la B-Box o derivada de ella, filtro y ajuste de nivel/salida.
\item Se incorporaron los valores actuales del canal SM-50: excitación de puente de \(\SI{10}{V}\), sensibilidad nominal de \(\SI{3}{mV/V}\), salida diferencial cruda de \(\SI{30}{mV}\), \(R_G=\SI{160}{\ohm}\), salida aproximada \(\SIrange{0}{10}{V}\) y eliminación de la resistencia serie de \(\SI{180}{\ohm}\).
\item Se actualizó el diagrama de bloques final para mostrar los dos canales con protección, amplificación instrumental, filtro pasabajo, ajuste de nivel/salida y conexión hacia BoomBox mediante FTP Cat5/8P8C.
\item Se incorporaron referencias completas guardadas en la carpeta del proyecto, incluyendo el manual de BoomBox, hoja de datos del SM-50, notas de aplicación de puentes Wheatstone, documentación de ruido, guías de KiCad y el texto de Pallás-Areny como referencia teórica principal.
\item Se generó una versión actualizada de la memoria en PDF y DOCX, junto con archivos auxiliares de cálculos y referencias.
\end{itemize}

\textbf{Resultados:}
La documentación quedó actualizada al estado actual del proyecto, con valores vigentes, referencias completas, diagrama final y separación explícita entre memoria final y memorias antiguas.

\textbf{Dificultades:}
La dificultad principal fue mantener consistencia entre memoria, diagrama, cálculos, referencias y bitácora, debido a que el proyecto tuvo varios cambios de valores durante su desarrollo.

% -------------------------
\subsection{Estado del proyecto al cierre de la actualización}
\textbf{Fecha:} 27/06/2026 \quad \textbf{Hora:} 06:00--07:00

\textbf{Objetivo:}
Registrar el estado final del proyecto al día de la actualización de la bitácora.

\textbf{Desarrollo:}
\begin{itemize}
\item El diseño queda definido como una PCB analógica de acondicionamiento para un sensor de torque SM-50 y sensores PT100, con bloques replicables para ampliar canales si se requiere.
\item La PCB recibe alimentación desde la BoomBox/B-Box por medio del conector 8P8C, usando \(\pm\SI{15}{V}\) con límite de \(\SI{100}{mA}\).
\item La adquisición final se realiza en la BoomBox; la PCB no digitaliza la señal, sino que la protege, excita, amplifica, filtra y ajusta antes de enviarla como señal analógica.
\item La interfaz física final queda definida mediante cable FTP Cat5, conectores 8P8C blindados y modo de entrada de alta impedancia diferencial.
\item El canal SM-50 queda como el canal con mayor exigencia de ganancia por partir de una señal diferencial de milivolts.
\item El canal PT100 queda definido con menor corriente de excitación para reducir autocalentamiento, manteniendo compensación por tres hilos y ajuste de offset.
\item La bitácora queda actualizada hasta el 27/06/2026 con las decisiones finales, manteniendo los documentos viejos solo como parte del proceso.
\end{itemize}

\textbf{Resultados:}
Se cuenta con una bitácora actualizada, coherente con la memoria final y con los valores actuales del diseño.

\textbf{Dificultades:}
Documentar el cierre de forma clara sin borrar el proceso previo, ya que las decisiones antiguas también fueron importantes para llegar a la versión final.

\begin{thebibliography}{99}

\bibitem{InterfaceSM50}
Interface, Inc. \emph{SM S-Type Load Cell: Datasheet}. Información técnica del sensor tipo puente Wheatstone.

\bibitem{InterfaceWiring}
Interface, Inc. \emph{Electrical \& Wiring Diagram}. Información de cableado para sensores de carga y torque.

\bibitem{TIBridgeSensors}
Texas Instruments. \emph{Signal Conditioning Wheatstone Resistive Bridge Sensors}. Application Report SLOA034.

\bibitem{TIStrainGauge}
Texas Instruments. \emph{Single-Supply Strain Gauge Bridge Amplifier Circuit}. SBOA247A.

\bibitem{ADBridgeCircuits}
Analog Devices. \emph{AN-43: Bridge Circuits}. Nota de aplicación sobre puentes resistivos y acondicionamiento.

\bibitem{ADBridgeADC}
Analog Devices. \emph{AN-96: Delta Sigma ADC Bridge Measurement Techniques}. Nota de aplicación sobre medición de puentes y resolución.

\bibitem{RuidoPCB}
R. Cárdenas Dobson. \emph{Atenuación de ruidos en circuitos electrónicos}. Apuntes del curso de proyecto, Universidad de Magallanes.

\bibitem{KiCadGuide}
Guía de proyecto. \emph{Guía para importar símbolos, huellas y modelos 3D en KiCad}.

\bibitem{Imperix}
imperix. \emph{Documentación técnica de módulos analógicos y BoomBox}. Referencia para interfaz analógica diferencial y conectores tipo RJ45.


\bibitem{PallasAreny}
R. Pallás-Areny, O. Casas y R. Bragós. \emph{Sensores y acondicionadores de señal: Problemas resueltos}. Ediciones de la Universitat Politècnica de Catalunya / Alfaomega, Barcelona, España, 2007.

\bibitem{MeasurementComputing}
Measurement Computing Corporation. \emph{Signal Conditioning \& PC-Based Data Acquisition Handbook}. 3rd ed., Measurement Computing Corporation, 2012.

\bibitem{BoomBoxManual}
imperix ltd. \emph{BoomBox User Manual}. Manual técnico de entradas analógicas, alimentación de sensores, filtros, ADC y conectores ModuLink/RJ45.

\bibitem{INA826}
Texas Instruments. \emph{INA826 Precision Instrumentation Amplifier}. Datasheet.

\bibitem{INA828}
Texas Instruments. \emph{INA828 Precision Instrumentation Amplifier}. Datasheet.

\bibitem{OPA2188}
Texas Instruments. \emph{OPA2188 Zero-Drift Operational Amplifier}. Datasheet.

\bibitem{REF3012}
Texas Instruments. \emph{REF3012 Precision Voltage Reference}. Datasheet.

\bibitem{TPS7A49}
Texas Instruments. \emph{TPS7A4901 High-Voltage Low-Noise Linear Regulator}. Datasheet.

\bibitem{MemoriasProceso}
Documentos internos del proyecto. \emph{Memorias preliminares, avances, diagramas y cálculos históricos de la PCB}. Usados como trazabilidad del proceso y no como valores finales.

\end{thebibliography}


\end{document}